% SPDX-License-Identifier: Apache-2.0
%
% HDMI in TxHDL. Built by //docs:hdmi. Every listing is included from
% the tree, and every printout and figure is what the build made by
% running the example.
\documentclass[journal,9pt]{IEEEtran}

% Listings of Rust set by rustfmt at 80 columns: 5.6pt is what fits
% the frame of a column (issue 1061).
\newcommand{\listingsize}{\fontsize{5.6}{6.6}\selectfont}
% SPDX-License-Identifier: Apache-2.0
%
% The house preamble, shared by every document in this directory. Loaded
% after \documentclass, before \title. Keeping it in one file is what
% stops two articles drifting apart in font, listing style or figure
% style.
% Computer Modern. IEEEtran selects Times (ptm) by default, so the face has
% to be chosen explicitly.
%
% Latin Modern is the Computer Modern variety used here, and the choice is
% forced rather than aesthetic. Under T1 encoding, plain `cmr` has no Type 1
% outlines unless cm-super is installed, and it is not in the pinned
% distribution, so pdflatex falls back to EC bitmaps and every glyph in the
% PDF becomes a Type 3 font. That was measured: the first build produced a
% document with 10 Type 3 fonts and no Type 1. Latin Modern is Computer
% Modern redrawn with a full T1 repertoire and real .pfb outlines, so the
% same design comes out scalable.
\usepackage[T1]{fontenc}
\usepackage[utf8]{inputenc}
\usepackage{lmodern}
% microtype is not in the pinned TeX distribution. emergencystretch alone
% keeps the two-column measure from producing overfull lines.
\setlength{\emergencystretch}{3em}

\usepackage{amsmath}
\usepackage{amssymb}
\usepackage{amsthm}
\usepackage{booktabs}
\usepackage{array}
% enumitem is not in the pinned TeX distribution; the list spacing below
% does what \setlist would have done.
\usepackage{float}
\usepackage{xcolor}
\usepackage{listings}
\usepackage{tikz}
\usetikzlibrary{arrows.meta,positioning,fit,backgrounds,calc}
\usepackage[hidelinks,bookmarks=true]{hyperref}

% Numbered, definition-style box for a rule the reader is meant to apply.
\theoremstyle{definition}
\newtheorem{rulebox}{Rule}
\newtheorem{example}{Example}

% Unnumbered asides.
\theoremstyle{remark}
\newtheorem*{tip}{Tip}
\newtheorem*{pitfall}{Pitfall}

% Tighter lists, in place of enumitem's \setlist.
\setlength{\topsep}{2pt}
\setlength{\itemsep}{1pt}
\setlength{\parsep}{0pt}

% Code in the running text. The typewriter face under T1 ligates `<<`
% and `>>` into guillemets, so a shift printed as \code{<<} came out as
% a French quotation mark (issue 571). microtype's \DisableLigatures is
% not in the pinned distribution, but the primitive it rests on is
% pdfTeX's own: \pdfnoligatures strips every ligature from the font it
% is given, and the current font inside \texttt is the typewriter face
% at the size in use, so each size is stripped the first time \code
% sets it. Code wants no ligature at all: two quotes are two quotes.
\newcommand{\code}[1]{\texttt{\ifdefined\pdfnoligatures\pdfnoligatures\font\fi #1}}
\newcommand{\term}[1]{\emph{#1}}

% Syntax highlighting. Four roles, distinguishable in colour and still
% distinguishable in greyscale, because an IEEE article gets printed: the
% keyword is the only bold role, the comment is the only italic one, and the
% two remaining roles differ in lightness as well as in hue.
\definecolor{kwblue}{RGB}{20,60,150}
\definecolor{cmtgreen}{RGB}{90,120,90}
\definecolor{strred}{RGB}{150,50,40}
\definecolor{numgray}{gray}{0.45}
\definecolor{framegray}{gray}{0.70}
\definecolor{bgshade}{gray}{0.97}
% A darker fill for figure cells. The listing background has to stay very
% light to keep code readable; a figure cell has to be clearly shaded or the
% distinction it draws is invisible in print.
\definecolor{cellshade}{gray}{0.90}

% The listing size is the document's choice, because it depends on the
% column: a two-column article sets `\listingsize` to 6pt and keeps its
% listings within 70 columns, or to 5.6pt when it includes source that
% rustfmt sets at 80, which is what fits a column's frame (issue 1061);
% a one-column document keeps the body size and includes source files
% at 80. `//docs:frame_test` checks every listing line against its frame. Lines are never broken; a line that does not fit
% overflows the frame, which is visible, rather than wrapping, which is
% not.
\providecommand{\listingsize}{\scriptsize}

% `'` is deliberately not a string delimiter: the language uses it for loop
% labels, as Rust does, so treating it as a quote would swallow the rest of
% the line.
\lstdefinestyle{txhdl}{
  basicstyle=\ttfamily\listingsize,
  keywordstyle=\bfseries\color{kwblue},
  commentstyle=\itshape\color{cmtgreen},
  stringstyle=\color{strred},
  numberstyle=\tiny\color{numgray},
  morekeywords={mod,use,pub,const,type,struct,enum,trait,impl,fn,async,let,mut,
    match,if,else,loop,while,for,in,break,continue,return,as,await,
    module,bus,channel,transaction,protocol,pipeline,flow,seq,config,
    port,stage,pipe,instance,connect,bind,tag,domain,blackbox,foreign,
    property,role,signal,handler,true,false,
    sequence,interface,modport,implements,package,import,var,wire,func,proc,
    case,default,next,fork,branch,join,state,fsm},
  morecomment=[l]{//},
  morestring=[b]",
  showstringspaces=false,
  columns=fullflexible,
  keepspaces=true,
  breaklines=false,
  % Framed, so a listing is bounded rather than floating in the column.
  frame=single,
  rulecolor=\color{framegray},
  framesep=3pt,
  backgroundcolor=\color{bgshade},
  xleftmargin=3pt,
  xrightmargin=3pt,
  aboveskip=6pt,
  belowskip=6pt,
  % Every listing is numbered and captioned, so it can be referred to by
  % number. `caption` without `float` numbers the listing and leaves it
  % where it was written, which is what a listing in running prose needs.
  captionpos=b,
  abovecaptionskip=2pt,
  belowcaptionskip=6pt,
}

% Figures drawn as text. Framed the same way, and with the highlighting
% switched off, because a box-and-arrow diagram is not code and half its
% words turning blue reads as an accident. Setting a style adds keys rather
% than clearing the ones already set, so the three roles are neutralised by
% hand rather than by leaving them out. Program output is printed in this
% style too, and a netlist's assignment can run past the frame, so a line
% that does not fit is broken and the rest indented; source never is.
\lstdefinestyle{ascii}{
  basicstyle=\ttfamily\listingsize,
  keywordstyle=\ttfamily,
  commentstyle=\ttfamily,
  stringstyle=\ttfamily,
  columns=fullflexible,
  keepspaces=true,
  breaklines=true,
  breakindent=2em,
  frame=single,
  rulecolor=\color{framegray},
  framesep=3pt,
  backgroundcolor=\color{bgshade},
  xleftmargin=3pt,
  xrightmargin=3pt,
  aboveskip=6pt,
  belowskip=6pt,
}
\lstset{style=txhdl}

% House TikZ styles. `shorten` lives on the arrow styles rather than on each
% arrow, so every arrow in the document gets the gap, including ones added
% later. TikZ clips a path at a node's border, so without it an arrowhead
% butts against the box with no gap at all. Keep `node distance` well above
% twice the shorten, or the arrow shrinks to nothing.
\tikzset{
  box/.style={draw,rounded corners=2pt,align=center,inner sep=4pt,
              font=\footnotesize},
  boxf/.style={box,fill=bgshade},
  lbl/.style={font=\scriptsize\itshape,align=center},
  fl/.style={-{Stealth[length=4pt]},shorten >=2.5pt,shorten <=2.5pt},
  fld/.style={fl,dashed},
}


% The contents list: the class gives a section's number 2.75em, which
% a Roman numeral past thirty overruns onto the title, so the box is
% wider here, and a subsection's entry starts where a title does.
\makeatletter
\def\l@section#1#2{\addpenalty{\@secpenalty}\addvspace{1.0em plus 1pt}%
    \@tempdima 5.75em \begingroup \parindent \z@ \rightskip \@pnumwidth%
    \parfillskip-\@pnumwidth {\bfseries\leavevmode #1}\hfil\hbox to\@pnumwidth{\hss #2}\par%
    \endgroup}
\def\l@subsection{\@dottedtocline{2}{5.75em}{5em}}
\def\l@subsubsection{\@dottedtocline{3}{10.75em}{5.5em}}
\makeatother

% SPDX-License-Identifier: Apache-2.0
% The boards' memory maps, written by //tools/memmap from the maps the
% designs decode (issue 444). A document asks for a table with
% \mmget{<what>}{<board>}, and a table the generator does not write
% stops the document rather than leaving a gap.
\input{tools/memmap/mm_tables}
\makeatletter
\newcommand{\mmget}[2]{%
  \@ifundefined{mm@#1@#2}%
    {\PackageError{memmap}{//tools/memmap writes no #1 for #2}%
      {Add it to tools/memmap/main.rs.}}%
    {\csname mm@#1@#2\endcsname}}
\makeatother


\InputIfFileExists{buildstamp.tex}{}{}
\providecommand{\buildstamp}{Draft build (outside the Bazel flow).}

% A range of the HDMI part, by its begin/end markers.
\newcommand{\hdmipart}[3]{%
\lstinputlisting[rangebeginprefix=//\ begin\{,rangebeginsuffix=\},%
rangeendprefix=//\ end\{,rangeendsuffix=\},includerangemarker=false,%
linerange=#1-#1,caption={#2},label={#3}]{lib/parts/src/hdmi.rs}}

\title{HDMI in TxHDL}

\author{Filip Filmar and Dragiša Janković%
\thanks{Both authors are independent. Filip Filmar is the
corresponding author, at f@hdlfactory.com; Dragiša Janković is
reached at linkedin.com/in/dragisa-jankovic.}%
\thanks{This is the tutorial on the HDMI part, for a reader who knows
the AXI link and its AXI-Lite bridge~\cite{axi}. Every listing is
included from the project repository \code{HDL/txhdl} at build time,
and every printout and figure is what the build produced by running
the example. The cover~\cite{cover} lists every document.}%
\thanks{The authors used a large language model, Claude, as an
assistant in exploring the concepts and in writing and constructing
the documents and the programs; every commit of the repository says
so and preserves the prompts that led to it, verbatim.}%
\thanks{\buildstamp}}

\markboth{HDMI in TxHDL}{Filmar and Janković: HDMI in TxHDL}

\begin{document}
\maketitle

\begin{abstract}
The Alinx AX7A200B development board transmits High-Definition Multimedia
Interface (HDMI) video through an onboard Silicon Image SiI9134
transmitter. The transceiver chip handles transition-minimized
differential signaling (TMDS) encoding and serialization. The FPGA
supplies a synchronous pixel clock, 24-bit parallel RGB color data,
horizontal and vertical synchronization strobes, and a display data enable
signal, initializing device registers over an I2C management bus. The
TxHDL implementation comprises two core units expressed within the lowered
subset: a video controller that sequences raster timing and scans a 12-bit
framebuffer accessible over AXI-Lite, and a dedicated I2C master that
releases transmitter reset and streams configuration tables. The build
suite exercises both units in a $16 \times 12$ scaled test mode against an
emulated AXI-Lite host and an I2C device model, validating pixel timing and
simulating synthesized netlists under \code{nvc} and Verilator. A
companion board demonstration configures standard $640 \times 480$
resolution at 60~Hz, synthesizes to a verified bitstream in Vivado, and
displays test patterns on physical HDMI monitors.
\end{abstract}

\section{What Is Here}
\label{sec:h-what}

The \code{hdmi} module in \code{txhdl\_parts} provides two hardware units.

\code{Hdmi} serves as the primary video peripheral. Running synchronously
on the pixel clock, it drives the transmitter inputs: 24-bit color bus
\code{rgb} (with red occupying the most significant byte), active-low
synchronization strobes \code{hsync} and \code{vsync}, and data enable
\code{de}, asserted during active video scan. Its internal memory stores
12-bit pixels (four bits each for red, green, and blue), with each
framebuffer entry replicated across a square of \code{1 << SHIFT} screen
pixels. Host processors update framebuffer contents across the unit's
AXI-Lite slave interface.

\code{I2cInit} automates transceiver initialization. It asserts external
chip reset, observes a stabilization delay upon reset deassertion, and
streams a sequence of configuration registers over I2C. The state machine
asserts \code{done} once the configuration sequence concludes, or asserts
\code{failed} if the transmitter fails to acknowledge any transmitted byte.

\section{The Raster}
\label{sec:h-raster}

A video frame comprises a rectangular raster of columns and rows divided
into active display intervals and blanking intervals. Each axis sequences
through active pixels, a front porch, a synchronization pulse, and a back
porch. For $640 \times 480$ resolution at 60~Hz, total line length spans 800
pixel clocks and frame height spans 525 scanlines. At a nominal 25.175~MHz
clock, pixel periods are 39.7~ns, with both synchronization pulses active
low. Listing~\ref{lst:h-modes} defines these timing parameters. The unit
parameterizes all eight values as type constants, enabling lightweight
test modes in simulation.

\hdmipart{modes}{640 by 480 at 60~Hz.}{lst:h-modes}

The controller tracks coordinates via column counter \code{hc} and row
counter \code{vc}. Combinational functions in Listing~\ref{lst:h-fns}
evaluate these counters to establish active display bounds and sync pulse
assertions. Synchronous block RAM reads sample the framebuffer into register
\code{px} on the rising clock edge. Synchronization strobes and data
enables register concurrently on the same edge, aligning timing controls
precisely with pixel data at the transmitter pins.

Three geometric helper functions operate across both axes generically
because compiler lowering binds constant function parameters by positional
argument order~\cite{i127}. Earlier compiler revisions bound parameters
by lexical identifier, forcing functions to duplicate names from
surrounding modules and requiring redundant axis-specific implementations.
Dedicated termination functions handle terminal column and row indices,
matching respective 8-bit and 7-bit counter bit widths.

\hdmipart{fns}{The raster's functions.}{lst:h-fns}

\section{The Framebuffer and the Bus}
\label{sec:h-bus}

The internal framebuffer provides $2^{15}$ addressable entries, indexed
with the row coordinate in the upper seven bits and the column in the lower
eight bits. Configured for $640 \times 480$ with a pixel scaling factor
\code{SHIFT} of 2, the active display spans $160 \times 120$ distinct
framebuffer pixels, leaving remaining storage unmapped.

Listing~\ref{lst:h-regs} declares the three-word register interface:
read-only register \code{status}, exposing vertical blanking state and
total frame count; read-write register \code{cursor}, holding active column
and row coordinates; and write-only register \code{pixel}, accepting pixel
color values. The \code{regmap!} macro synthesizes read multiplexers, write
strobes, and host software field accessors from this unified description.

\hdmipart{regs}{The video peripheral's words.}{lst:h-regs}

Writing to \code{pixel} stores color data at the current cursor coordinates
and auto-increments the pointer. Upon crossing a scanline boundary, the
cursor wraps to the start of the subsequent row, returning to the origin
after completing the final line. Software sets the cursor once and streams
continuous pixel spans via fixed-burst AXI transfers, which the bridge
translates into consecutive single-beat writes at the constant register
address.

\hdmipart{state}{The video peripheral's state.}{lst:h-state}

\hdmipart{run}{The video peripheral's step.}{lst:h-run}

Output color expansion constructs intermediate signals for each color
channel before bitwise concatenation. Earlier lowering passes failed to
resolve chained expressions following function calls, incorrectly treating
them as static constants (issue~128)~\cite{i128}.

\section{Configuring the Chip}
\label{sec:h-i2c}

The I2C physical bus comprises two open-drain signals, serial clock (SCL)
and serial data (SDA), pulled up to positive supply rails by onboard
resistors. Bus masters initiate transfers with a START condition (SDA
transitioning low while SCL remains high) and terminate with a STOP
condition (SDA transitioning high while SCL remains high). Data octets
transfer most-significant bit first, shifting while SCL is low and
sampling while SCL is high, followed by an acknowledge cycle during which
the recipient pulls SDA low.

The I2C master subdivides each bit period into four quarters of \code{DIV}
clock cycles, driving SCL low during the first and fourth quarters. Each
transaction comprises 29 four-quarter intervals: START framing, 7-bit slave
address plus write flag with ACK, register index with ACK, payload data with
ACK, and STOP framing. Control outputs \code{scl\_low} and \code{sda\_low}
pull bus lines low actively. Inbound signal \code{sda\_in} allows the master
to sample ACK status at the midpoint of the third quarter.

The controller partitions into two concurrent processes. The timing
process maintains a divider counter modulo \code{DIV}, generating
quarter-bit phase strobes and registering external status signals
\code{done} and \code{failed}. The sequencing process expresses transfer
protocols directly: asserting reset for \code{HOLD} intervals, awaiting
post-reset recovery, and iterating through configuration entries. Each
table entry executes START generation, three serialized bytes with
acknowledge checks, and STOP generation. Compiler lowering assigns these
sequential waits into state registers and synthesizes loop counters
automatically, eliminating manual phase and step tracking. The sequence
terminates in a permanent wait state once \code{written} asserts.

\hdmipart{i2cstate}{The master's registers.}{lst:h-i2cstate}

\hdmipart{i2c}{The master's two processes.}{lst:h-i2c}

Listing~\ref{lst:h-table} provides the configuration table programmed by
the master: enabling operational power states for the board-wired input
buses and setting output format to DVI mode without auxiliary HDMI
info-packets. Because reference documentation omitted register
specifications, these values were derived and later verified on physical
hardware, where the SiI9134 accepted each transaction and activated video
output. Helper functions in the source supply table entries to the sequencer.

\hdmipart{table}{The chip's configuration, confirmed on the
board.}{lst:h-table}

\section{The Run}
\label{sec:h-run}

Example \code{ex\_hdmi} instantiates the video controller behind
\code{LiteBridge<1, ..>} and a host transaction tracker in a scaled mode
of $16 \times 12$ active pixels. Framebuffer pixels are scaled $2 \times 2$,
yielding an $8 \times 6$ backing store. The host client resets the cursor
to the origin and paints all 48 pixels across a single fixed burst, ramping
red horizontally, green vertically, and blue diagonally. The host then polls
\code{status} until two subsequent video frames display.

Concurrently, the I2C master executes against \code{I2cDevice}, a bus model
that acknowledges received octets and logs transaction data. Test
assertions verify that transmitted values match the configuration table.

The verification harness reconstructs full frames from active pixels
marked by \code{de} following vertical sync, comparing all 192 screen
pixels of the final frame against expected pattern colors. A coordinate
shift of a single row or column would immediately trigger mismatch
assertions at $2 \times 2$ pixel borders.

\begin{figure*}[t]
\centering
\input{hdmi_timing.tex}
\caption{The run: the pixels written on \code{w}, the syncs and the
enable frame after frame, and the chip's reset and
the I2C lines while the two transactions go out, until \code{done}.}
\label{fig:h-run}
\end{figure*}

\lstinputlisting[caption={\code{ex\_hdmi.rs}. Run with
\code{bazel run //lib/examples:ex\_hdmi}.},%
label={lst:h-ex}]{lib/examples/ex_hdmi.rs}

\lstinputlisting[style=ascii,caption={What \code{ex\_hdmi} printed when
the build ran it: the two transactions the device was sent, and the
last whole frame, one screen pixel per framebuffer pixel.},%
label={lst:h-out}]{out_hdmi.txt}

Both units undergo lowering, and the build suite simulates generated
netlists at their port boundaries under \code{nvc} and Verilator against
this test stimulus. Unit tests within the module verify sync pulse timing,
data enable delays, and I2C protocol waveforms against reference mode
definitions.

\section{On the Board}
\label{sec:h-board}

The board-level demonstration is implemented in \code{//hdmi} without an
active bus host. Netlist generation initializes the framebuffer with a
static test image comprising eight vertical color bars, a grayscale ramp,
and a bounding border, with AXI-Lite inputs tied inactive. On the flagship
system, the processor core accesses the peripheral at slot~3 of its
peripheral map (\mmget{at}{hdmi}).

\lstinputlisting[caption={\code{hdmi/src/netlist.rs}: the netlist, with
the test picture.},label={lst:h-netlist}]{hdmi/src/netlist.rs}

An MMCM synthesizes the pixel clock from the 200~MHz board reference:
dividing by 10 and multiplying by 63 establishes a 1260~MHz internal
oscillator, which divided by 50 produces 25.2~MHz (within 0.1\% of the
25.175~MHz nominal rate). An \code{ODDR} primitive inverts the clock driving
the transceiver, centering transmitter sampling edges within pixel data
eyes. I2C pins configure as open-drain buffers driven low when pulling and
tri-stated otherwise. The clock manager lock signal provides system reset,
passing through dual synchronizers clocked by the pixel domain. Lowered
modules reset exclusively via their \code{rst} input (issue~509).

\lstinputlisting[caption={\code{hdmi/board/hdmi\_demo.v}: the
demonstration.},label={lst:h-demo}]{hdmi/board/hdmi_demo.v}

Vivado synthesizes the lowered units and top wrapper via
\code{bazel build //hdmi:demo\_synth} without errors or critical warnings;
the 87 informational warnings reflect unused tied-low AXI-Lite inputs.
Implementation via \code{bazel build //hdmi:demo\_pnr} achieves timing
closure with 35.872~ns setup slack and 0.164~ns hold slack across 398
endpoints. The design utilizes 60 LUTs, 86 flip-flops, 12 block RAM tiles,
and 36 I/O pins. Because external PCB trace delays are unmodeled in
constraints, reported margins apply to internal FPGA logic.

Board timing and LED pins reside in \code{hdmi/board/ax7a200\_hdmi.xdc},
while the 32 transceiver pin constraints are defined in
\code{hdmi/board/hdmi\_pins.xdc}, which the flagship build imports directly.

\lstinputlisting[caption={\code{hdmi/board/ax7a200\_hdmi.xdc}: the
board's pins.},label={lst:h-xdc}]{hdmi/board/ax7a200_hdmi.xdc}

\lstinputlisting[caption={\code{hdmi/board/hdmi\_pins.xdc}: the
chip's pins.},label={lst:h-pins}]{hdmi/board/hdmi_pins.xdc}

\section{What Is Not Done}
\label{sec:h-not}

The hardware demonstration was validated on physical hardware. On September
19, 2026, an HDMI monitor attached to the AX7A200B displayed the test
pattern stored in the framebuffer: eight color bars (white, yellow, cyan,
green, magenta, red, blue, and black) spanning the upper two-thirds, a
grayscale gradient below, and a white border. Diagnostic LEDs confirmed
clock lock and configuration success, with every I2C transaction
acknowledged.

These results confirmed three initial design parameters. First, the
configuration register sequence succeeded, activating TMDS video output.
Second, 24-bit RGB wiring matches transceiver pin ordering, as evidenced
by correct color bar sequence (any red-blue swap would transform yellow into
cyan). Third, 3.3~V LVCMOS signaling operates correctly, aligning with
vendor reference designs.

Board validation identified one revision difference. Vendor reference
designs drive two reset pins on FPGA balls L18 and Y17 to support hardware
revisions 1.0 and 2.0. The design now drives both pins simultaneously to
guarantee reliable reset release across board revisions.

External pin timing constraints remain uncharacterized; while stable video
confirms adequate operational margin, precise setup and hold windows at the
connector are unmeasured.

The peripheral currently operates entirely within the pixel clock domain,
including its bus interface. Connecting a processor running on a separate
system clock requires clock-domain crossing across all five AXI-Lite
channels. In TxHDL, module \code{ChanCdc} provides channel-level
synchronization (issue 131)~\cite{i131}. The flagship SoC instantiates five
instances of handwritten module \code{chan\_cdc} to bridge these domains.

Current architectural boundaries fix the framebuffer at 12 bits per pixel
with maximum dimensions of $256 \times 128$. Video timing constants are
fixed at compile time by type parameters, and I2C initialization executes
a single static table.

\IEEEtriggeratref{1}
\begin{thebibliography}{9}

\bibitem{axi}
F.~Filmar and D.~Janković, ``AXI in TxHDL,'' project repository
\code{HDL/txhdl}, \code{//docs:axi}, 2026.

\bibitem{i127}
Issue 127, ``fix(lower): a function's const parameters are bound by
name, not by position,'' project repository \code{HDL/txhdl}, 2026.

\bibitem{i128}
Issue 128, ``fix(lower): a method chain on a function call becomes a
constant naming missing variables,'' project repository
\code{HDL/txhdl}, 2026.

\bibitem{i131}
Issue 131, ``feat(test): co-simulate a unit of several clocks, and
lower a clock-crossing channel,'' project repository \code{HDL/txhdl},
2026.

\bibitem{cover}
F.~Filmar and D.~Janković, ``TxHDL documents,'' project repository
\code{HDL/txhdl}, \code{//docs:cover}, 2026.

\end{thebibliography}

\end{document}
